\section{Lean formalization: definitions, targets, and the completion path}
\label{sec:lean-plan}

The formalization has three distinct endpoints. First, executable native
definitions fix the finite objects and prevent an informal theorem from
changing shape during implementation. Second, native correctness lemmas
establish the finite mathematics in Lean. Third, embedded derivations
establish the same assertions in the specified iPRA calculus. This report
provides a small compiled part of the second endpoint, compiled definitions
for several of the remaining targets, and a detailed plan for the third.

\subsection{The existing repository interfaces}

The pinned repository's
\path{ConciseIPRA/Lean/ConciseIPRA/PR.lean} defines an inductive type
\texttt{PR k} with constructors for zero, successor, projection,
composition, and primitive recursion. It also defines \texttt{PTm n},
the terms over $n$ free variables. The syntax file supplies
\texttt{Fm PTm n}; the deduction file supplies the derivation type
\texttt{Prf Ax Gamma phi} for intuitionistic natural deduction with
equality. The arithmetic axiom family is \texttt{PRAAx}.

In particular, the repository's induction constructor accepts a proof
that its formula is quantifier-free. The derived rule
\texttt{PRA.qfInd} takes the base derivation and a derivation of the
successor step under the shifted context. A completed theorem in the
present report should end with a declaration of the following shape:
\begin{verbatim}
def target : Fm PTm parameterCount := ...
theorem target_derivation : Prf PRAAx [] target := ...
\end{verbatim}
The ellipses here denote implementation work explicitly listed below;
they are not part of a claimed compiled theorem.

The module \texttt{NativeBiInterpretation.lean} supplies
\texttt{pra\_iff\_native}. It transports an already supplied iPRA
derivation to the native presentation, using \texttt{NFm.ofPRA}.
It does not turn an arbitrary Lean proposition into an iPRA theorem.
The archive's \texttt{ObjectInterface.lean} records these actual types,
fixed PR name symbols, syntactic implication composition, and the
transport step. That integration file has not been compiled here,
because it requires the full repository library and Mathlib. Its status
is separate from the three standalone compiled modules.

\subsection{The compiled native definitions}

\paragraph{Finite search and contracts.}
\texttt{FiniteCore.lean} uses only Lean's standard distribution. It defines
a bounded scan \texttt{firstHit p n} that retains the first successful
index in $0,\ldots,n$. If no success exists it returns $n$; therefore
success is always tested by applying \texttt{p} to the returned value,
not inferred from the index alone. This convention differs from the
sentinel implementation in Lemma~\ref{lem:bounded-choice}, but both
are explicitly bounded and their success conditions are unambiguous.

The file proves the bound on the returned index, completeness from a
supplied successful index, and the failure certificate for every
inspected index. Its \texttt{energy\_firstHit} lemma proves success
whenever a natural-valued energy starts below $B$ and drops by at least
one on each failed test through $B$. The proof maintains
$E(i)+i\le B$ under a hypothetical all-failure prefix. At $i=B+1$
that is impossible. The file also defines disjoint endpoints and
proves their monotonicity, composes finite contracts, and packages the
selected-window bound.

\paragraph{Polynomial and observation targets.}
\texttt{Targets.lean} fixes concrete native data:
\begin{itemize}
\item rational polynomials as coefficient lists, with coefficient lookup,
evaluation, addition, multiplication, degree, nonzero tests, and weighted
finite coefficient sums;
\item rational point lists, exact pairwise distinctness, the common-zero
test for $y^2=f(x)$ and $A(x)+B(x)y=0$, and the norm polynomial;
\item finite lists of bivariate terms, coefficient collection, substitution,
and polynomial iteration truncated modulo $X^{K+1}$;
\item regular rational-name predicates, finite regularity tests, finite
variation, and finite window-stability tests;
\item bivariate spanning-certificate data, matrix rows, ideal-multiplier
witnesses, dimension checks, exact border-identity checks, and the
common-zero predicate.
\end{itemize}
Definitions named \path{HyperellipticFiberTarget},
\path{ImplicitCoefficientStabilityTarget},
\path{FiniteVariationTarget}, and \path{FiniteAlgebraTarget}
are propositions specifying the remaining theorems. They are \emph{not}
theorems or assumed axioms. The bivariate certificate target is a concrete
two-variable instance of the general $r$-variable theorem in the paper.
Its checker verifies that the basis is nonempty and begins with $1$,
that all matrix and witness rows have the correct lengths, and that
every required polynomial identity holds after coefficient collection.

The finite-variation target uses the sufficient hypothesis
\[
 \operatorname{variation}(T)<(B+1)\epsilon
\]
to conclude that one of $B+1$ windows succeeds. This matches the
floor-budget construction in Section~\ref{sec:stability}. The native
functions in these statements can be arbitrary host functions. Their
additional fixed-PR representation is required when constructing the
object-theory schema; the native type alone does not certify it.

The source convention for the degree of the zero polynomial is zero.
This is harmless in the fiber target because its hypotheses imply
that the norm polynomial is nonzero. Polynomial multiplication may
retain trailing zero coefficients; coefficient comparison and the
implemented degree ignore those zeros. Weighted-norm lemmas must assume
their weight is nonnegative, and the analytic lemmas in the paper use
a strictly positive rational weight.

\paragraph{The Gaussian-rational branch prototype.}
\texttt{FiniteBranch.lean} adds the paired-coordinate version of the
polynomial algorithms and records the finite analytic goals in a
\texttt{KernelGoals} structure. No instance of that structure is supplied.
It also checks five individual coefficients of the branch of $Y^2=1+X$
by kernel computation. The companion \texttt{ANALYTIC\_PLAN.md}
specifies the finite norm, perturbation, and name proofs still required.

\subsection{Constructing the embedded arithmetic layer}

The following sequence gives concrete milestones. Each milestone has
a finite correctness statement and an indicated proof method.

\paragraph{1. Natural coding and arithmetic.}
Choose one total PR pairing function with bounded projections and a
finite-list code with PR length and entry functions. A radix list code
with an explicitly stored base is sufficient whenever a finite maximum
of the entries is available; a standard nested pairing is sufficient
otherwise. Integers may be represented by pairs of naturals, with
equality $(a,b)\simeq(c,d)$ defined by $a+d=b+c$. Represent rationals by
an integer numerator and a positive denominator. Do not assume that
arbitrary host \texttt{Rat} computations are already object-theory
functions.

Build the PR definitions and derivations for arithmetic, comparison,
bounded sums, bounded maxima, and list operations. Prove validity
preservation separately. Cross-multiplication turns the rational
inequalities needed below into integer inequalities and then PR Boolean
tests. No quotient construction is needed inside iPRA.

\paragraph{2. Bounded formulas and proof construction.}
Implement the formula-to-Boolean translation of
Lemma~\ref{lem:bounded-booleanization} as a metalevel transformation
on a fixed formula syntax, emitting both a PR definition and its
equivalence derivation. Implement the bounded search derivations next.
The syntactic compiler then uses substitution, intuitionistic logical
rules, and \texttt{PRA.qfInd} on the emitted Boolean invariant. In
particular, it must never invoke ordinary Lean induction on an arbitrary
semantic proposition and declare the resulting proof to be an embedded
derivation.

\paragraph{3. Finite algebra.}
Prove coefficient-array convolution, distributivity, the difference of
powers identity, and synthetic division. The normal-form soundness
argument in Proposition~\ref{prop:polynomial-normalization} gives a
reusable source of finite identity derivations. Its ring operations are
fixed PR definitions with supplied arithmetic laws. Parameterized matrix
identities still need their uniform generators or direct proofs.

\paragraph{4. Analytic finite kernels and names.}
Use paired rational coordinates for Gaussian rationals. Prove their
$\ell^1$ triangle and product bounds. Prove (A1)--(A9) on finite
polynomial arrays; the induction invariants are the weighted norm bound,
the bounded list of zero low coefficients, and the rational perturbation
inequality. Fix the requested coefficient degree and both precision
indices before induction. Then instantiate these finite lemmas at the
explicit indices (A12)--(A17). The input regularity and root equation are
universal arithmetic assumptions that are used at the finite numerical
instances demanded by a particular conclusion.

For represented equality, adopt one fixed convention. For example, two
regular names $a,b$ represent the same value when for every requested
$n$ their sufficiently reindexed approximants differ by at most
$2^{-n}$. State the exact reindexing in each operation and verify its
transitivity by a further fixed precision shift. No equality decision
procedure is introduced. For root uniqueness, use the finite
contraction-with-error inequality and choose its error below the
requested tolerance.

\paragraph{5. Finite stability.}
Prove telescoping and the triangle inequality for a supplied finite
sequence, then the disjoint-increment sum bound. A failed window
contains a boundedly found witness pair. The finite assignment table
records which component pays for each failed window. Its invariant is
that all assigned indices lie in the processed prefix, that the classes
partition that prefix, and that each count respects its proved budget.
This produces the successful index needed by the already compiled
native search theorem, and separately produces an embedded proof of
the coded selector's correctness.

\paragraph{6. Hyperelliptic and finite algebra counts.}
For the norm theorem, carry the point groups, the removed product, and
the residual polynomial. The central invariant is $N=CQ$ together with
the degree sum and the number of points accounted for. The multiplicity
two lemma is finite rational field arithmetic. For the certificate
theorem, carry reduction vectors and the evaluation matrix of the
separator polynomials on the supplied point list. Bounded rational row
reduction tracks an invertible companion matrix, as described in
Section~\ref{sec:repository-repair}. These invariants are Boolean
equalities over finite codes, so their induction is permitted.

\paragraph{7. Assembly and audit.}
Construct the concluding \texttt{Prf PRAAx} terms. Audit every premise:
the polynomial identities, root equation, derivative gap, continuity
modulus, finite variation bound, spanning identities, or finite covering
theorem must either be input data with the stated contract or have an
actual preceding derivation. Only after this assembly invoke
\texttt{pra\_iff\_native}. There is no need to reprove bi-interpretability
for every application.

\subsection{What the formalization does and does not certify}

The lemmas in \texttt{FiniteCore.lean} prove native finite-search and contract
facts. Their printed axiom dependencies include Lean's standard
\texttt{propext} and \texttt{Quot.sound} for the arithmetic tactic
proofs; the contract-composition proof has no axiom dependencies.
They contain no \texttt{sorry}, no new axioms, and no use of
\texttt{Classical.choice}. This does not by itself certify that these
native proof terms belong to iPRA. The explicit arithmetic invariants
and the planned construction of \texttt{Prf PRAAx} terms are what address
that additional requirement.

The five computation proofs in \texttt{FiniteBranch.lean} have a
different printed dependency set: it includes \texttt{Classical.choice}
as well as \texttt{propext} and \texttt{Quot.sound} through their
native rational/container infrastructure. They are disclosed native
metatheory computations. The intended object proof instead uses the
raw integer/denominator representation and finite arithmetic arguments
given in the report; these computation proofs are not imported as
unexamined iPRA axioms.

The report's human proofs are substantially more developed than the
current Lean implementation. The remaining work is finite and specified
for the proved lemmas, but it includes real implementation and proof
construction: no blanket promise that a future \texttt{ring} or
\texttt{simp} call will discharge it is being made. The general
Diophantine approximation and height theorems identified in the audit
are further mathematical problems, not merely this implementation work.
